\documentclass[12pt]{article}
\usepackage[a4paper,margin=1in]{geometry}\usepackage[T1]{fontenc}
\usepackage{lmodern,microtype}\usepackage{amsmath,amssymb,amsthm,mathtools}
\usepackage{booktabs,array,enumitem,xcolor}\usepackage[numbers,sort&compress]{natbib}
\usepackage[colorlinks=true,linkcolor=blue!55!black,citecolor=blue!55!black,urlcolor=blue!55!black]{hyperref}
\linespread{1.07}
\newtheorem{theorem}{Theorem}[section]\newtheorem{proposition}[theorem]{Proposition}
\theoremstyle{remark}\newtheorem{remark}[theorem]{Remark}
\newcommand{\norm}[1]{\left\lVert#1\right\rVert}

\title{Expanding the Modulus-Selected Edge Cloud Does Not Repair Transport\\
\large A Rank-Two-to-Rank-Thirty-Two Principal-Angle Audit}
\author{Liang Wang}\date{July 2026}
\begin{document}\maketitle
\begin{abstract}
The outer quartet may fail Haar transport because four modes are too few to
absorb rotating spectral directions.  We test this natural repair by
selecting the $k$ largest-modulus modes at both adjacent levels for
$k=2,4,6,8,12,16,24,32$ and comparing the embedded right and left invariant
spaces.

The repair fails on the declared rank grid.  None of the 32
rank/transition/channel cases has both right and left maximum principal
sines below $0.5$.  The best joint sine is $0.69373$.  On the
$0.04\to0.02$ cases rank four is optimal among the tested ranks; on the
$0.02\to0.01$ cases rank two is optimal.  For ranks above four, the joint
maximum sine often approaches one and reaches $0.999999993$.

We also give an exact finite example showing that nested cloud dimensions do
not force monotonic improvement of the largest principal angle: adding one
orthogonal direction can change a zero angle to $\pi/2$.  The result rejects
only equal-rank, top-modulus clouds under the naive Haar map.  It does not
rule out adaptive clusters, unequal-rank transport, a renormalized embedding,
or a scalar divisor limit.
\end{abstract}

\section{The cloud-enlargement hypothesis}

The fixed quartet is necessarily only a local block: an infinite spectral
count requires packet rank to grow.  After RH-202 and RH-208 reject its naive
Haar transport, a natural hypothesis is that omitted neighboring modes carry
the missing directions.  If so, enlarging the selected edge cloud should
decrease the largest principal angles.

We test this hypothesis without changing two ingredients:
\begin{enumerate}[leftmargin=2em]
\item modes are selected by decreasing eigenvalue modulus;
\item the physical interlevel map is the dyadic Haar embedding.
\end{enumerate}
Only the cloud rank changes.

\section{Selection and metric}

Let $\lambda_1,\ldots,\lambda_n$ be ordered so that
\begin{equation}\label{eq:order}
 |\lambda_1|\ge\cdots\ge|\lambda_n|.
\end{equation}
Let $E_c(k)$ and $E_f(k)$ be the right invariant spaces spanned by the first
$k$ modes at adjacent levels, and let $F_c(k),F_f(k)$ be the corresponding
left spaces.

For orthonormal bases $Q_c,Q_f$, the singular values of $Q_f^*JQ_c$ are the
principal cosines \citep{StewartSun1990}.  Define
\begin{equation}\label{eq:joint}
 g_k=\max\{\sin\theta_{\max}(JE_c(k),E_f(k)),
 \sin\theta_{\max}(JF_c(k),F_f(k))\}.
\end{equation}
The predeclared exploratory gate is
\begin{equation}\label{eq:gate}
 g_k<0.5.
\end{equation}
Passing would not prove transport, but would identify a cluster worth a more
expensive oblique-projector and contour analysis.

\section{Why monotonicity is not automatic}

\begin{proposition}[Nesting does not improve the worst angle]\label{prop:nonmono}
There exist nested pairs $E_c(1)\subset E_c(2)$ and
$E_f(1)\subset E_f(2)$ such that the largest principal angle is zero at rank
one and $\pi/2$ at rank two.
\end{proposition}

\begin{proof}
In $\mathbb R^3$, take
$E_c(1)=E_f(1)=\operatorname{span}\{e_1\}$,
$E_c(2)=\operatorname{span}\{e_1,e_2\}$, and
$E_f(2)=\operatorname{span}\{e_1,e_3\}$.  The rank-one spaces coincide, but
the second principal cosine at rank two is zero.
\end{proof}

Thus an increasing mean overlap can coexist with a worsening maximum angle.
The latter is the relevant quantity for invertible packet correspondence.

\section{Thirty-two physical rank cases}

The joint maximum sines are:
\begin{center}
\scriptsize
\begin{tabular}{llrrrrrrrr}
\toprule
step/side & $k=2$ & $4$ & $6$ & $8$ & $12$ & $16$ & $24$ & $32$\\
\midrule
$0.04\to0.02$ left
& $.8108$ & $.7654$ & $.9769$ & $.9998$ & $.9996$ & $.9996$ & $1.0000$ & $1.0000$\\
$0.04\to0.02$ right
& $.8133$ & $.7669$ & $.9769$ & $.9989$ & $.9994$ & $.9999$ & $1.0000$ & $1.0000$\\
$0.02\to0.01$ left
& $.6948$ & $.8217$ & $.9719$ & $.9854$ & $.9982$ & $.9996$ & $1.0000$ & $1.0000$\\
$0.02\to0.01$ right
& $.6937$ & $.8239$ & $.9755$ & $.9870$ & $.9992$ & $.9999$ & $1.0000$ & $1.0000$\\
\bottomrule
\end{tabular}
\end{center}
No value passes \eqref{eq:gate}.  More importantly, ranks six and above are
uniformly worse than the best low-rank choice in every case.

\section{Best-rank pattern}

The rank distribution of the minimizers is
\begin{equation*}
 k=4\text{ in two cases},\qquad
 k=2\text{ in two cases},
\end{equation*}
with no minimizer at a larger rank.  The quartet remains the best tested
choice on the first transition, while one conjugate pair is best on the
second.

This pattern argues against a simple rule of the form ``take more outer
modes until transport closes.''  The extra modes are not merely a buffer;
some bring nearly orthogonal directions into the worst-angle metric.

\section{Interpretation of near-unit angles}

A maximum sine near one means that at least one direction in one selected
cloud is almost orthogonal to the other after embedding.  It does not imply
that all directions disagree: mean principal cosines can remain moderate.
However, one lost direction is enough to make an equal-rank transport map
ill conditioned or singular.

For oblique nonnormal packets the situation is stricter still.  Good right
angles would not guarantee good left angles, which is why the joint metric
in \eqref{eq:joint} is used.

\section{Worst-angle versus average-overlap diagnostics}

The mean principal cosine answers how much of a cloud overlaps on average;
the minimum cosine answers whether every direction is recoverable.  A
determinant factor of equal degree needs the latter: losing one direction can
change one root or make the transport singular even when most energy is
captured.  This explains why low-rank postblock energy compression from
earlier papers does not automatically imply spectral-cloud transport.

For an unequal-rank future map, the appropriate diagnostics change.  One
should record injection angles from the coarse packet into a larger fine
packet, the dimension of the captured image, and the residual fine
complement.  Such a test may succeed even when the equal-rank largest angle
is near $\pi/2$; it is outside the claim of RH-209.

\section{What has been ruled out}

The finite result rejects the compound rule
\begin{equation}\label{eq:rejected}
 \text{same rank at both levels}
 +\text{ largest-modulus selection}
 +\text{ Haar embedding}
\end{equation}
on the specified rank grid and four physical cases.

It does not reject:
\begin{enumerate}[leftmargin=2em]
\item a contour selected by branch continuation rather than modulus;
\item different coarse and fine cloud ranks with an injective map;
\item weighted subspace metrics adapted to source observability;
\item a scale-dependent coordinate dilation or other renormalized embedding;
\item convergence only of scalar characteristic data.
\end{enumerate}

\section{Predeclared alternatives for a later cloud audit}

Three nonmodulus rules are sufficiently concrete to test without post hoc
selection: continue contours from the RH-204 branch labels; select modes by
a joint modulus/residue threshold; or select the smallest conjugation-closed
cluster whose transfer moments meet a declared error gate.  Each rule must
freeze its threshold before inspecting transport angles.  Otherwise cloud
growth can be tuned to manufacture an apparent match.

\section{Consequence for Gate A}

The quartet should remain a local diagnostic and a branch-label seed, but it
should not be promoted to a nested raw eigenspace shell.  Likewise, blindly
growing a modulus cloud is not supported.  The strongest surviving signal is
the dual-channel scalar divisor coherence found in RH-207.

This motivates a route decision: determine whether convergence of spectral
divisors can be pursued independently of raw state transport.  RH-210 proves
that state convergence is sufficient but not necessary for divisor
stability and records the resulting Gate-A pivot.

\section{Claim boundary}

The nonmonotonicity proposition is exact.  The rank-grid obstruction is
finite and floating.  It is not an all-rank or all-level theorem and does
not exclude adaptive nonmodulus clouds.  No infinite determinant, counting
law, arithmetic trace, or Hilbert--P\'olya statement follows.

\bibliographystyle{plainnat}\bibliography{references}\end{document}
